\section{条件期望是利用信息的最佳预测}
\label{sec:05-Probability/04-Expectation-and-Moments/04}

你是一家电力公司的调度员。没有任何额外信息时，你对明天用电需求的最优预测就是历史总均值——假设为 500 MW（方差一节证明过：平方误差下，\(\E[X]\) 是最优常数预测）。现在气象局发来明天预报温度：35 度，高温。经验告诉你，高温天和凉爽天用电量完全不在一个级别。如果你继续用 500 MW 做预测，就是在主动忽视一条有用的信息。

问题变得具体：知道温度 \(Y\) 之后，给每个可能的温度值配上一个预测值——这套 ``随信息变化的预测函数'' 怎么写最合理？平方误差再次成为判据：寻找函数 \(g(Y)\)，极小化 \(\E[(X-g(Y))^2]\)。

\subsection{从离散分层出发}

最直观的入口是 \(Y\) 只取几个离散值。把数据按 \(Y\) 的值切成分层：\(Y=y_1\) 的那批样本算一个均值，\(Y=y_2\) 的算另一个。

算一下。设温度只有两种情形：高温天（概率 0.3），低温天（概率 0.7）。高温天的历史用电均值 \(\E[X\mid\text{高温}]=650\) MW，低温天 \(\E[X\mid\text{低温}]=436\) MW。你的预测函数自然改为

\[
g(\text{高温})=650,\qquad g(\text{低温})=436.
\]

这就是条件期望 \(\E[X\mid Y]\) 在离散情况下的形态。对每个固定的 \(Y=y\)：

\[
\E[X\mid Y=y]=\sum_x x\,P(X=x\mid Y=y),
\]

和普通期望的定义结构一模一样，只是把无条件概率换成了条件概率。把所有这些数收集起来，看成 \(y\) 的一个函数，再把随机的 \(Y\) 代进去：

\[
\E[X\mid Y]=g(Y),\quad\text{其中 }g(y)=\E[X\mid Y=y].
\]

\(\E[X\mid Y]\) 本身是一个随机变量——在你看到温度之前，你不知道预测值会是多少；看到了，它就坍缩成一个具体数字。

另一个例子帮你识别指示变量的模式。袋中先随机选袋（红袋概率 0.5，蓝袋 0.5），再从选中袋子中抽一次。红袋内成功概率 0.8，蓝袋内 0.3。令 \(X\) 为 ``是否成功'' 的指示变量，\(Y\) 为袋子颜色。则

\[
\E[X\mid Y=\text{红}]=0.8,\qquad
\E[X\mid Y=\text{蓝}]=0.3.
\]

因为指示变量的期望就是事件概率，条件期望 \(\E[\mathbf 1_A\mid Y]\) 正好就是 \(P(A\mid Y)\)——条件期望这个概念同时消化了条件概率。

\subsection{塔式性质：分层之后再合回来}

分层计算之后，可以验证它们和总体均值是否衔接得上。总期望等于每层的期望按层的概率加权平均：

\[
\E[X]=\sum_y \E[X\mid Y=y]\,P(Y=y)=\E\bigl[\E[X\mid Y]\bigr].
\]

用电例子：\(0.3\times650+0.7\times436=195+305.2=500.2\approx500\)（小数由四舍五入）。各层均值加权回来，恰好恢复总体均值。

这条性质在教材里叫全期望公式或塔式性质。名字 ``塔'' 来自更一般的形式：如果信息集合 \(\mathcal G\) 比 \(\mathcal H\) 更粗糙（即 \(\mathcal G\) 包含的信息更少），那么

\[
\E\bigl[\E[X\mid\mathcal H]\mid\mathcal G\bigr]=\E[X\mid\mathcal G].
\]

外层对 \(\mathcal G\) 取条件期望相当于故意忘掉一部分信息，因此和一开始只用粗信息预测结果相同。从外向里依次应用：先用较多信息做预测，再把这个预测按较少信息做平均，得到的结果等于直接用较少信息做的预测。

\subsection{全方差公式：总波动拆成两层}

上节全期望公式揭示了条件均值如何 ``回归'' 为总体均值。现在把方差的来源也拆开。条件期望给每个 \(Y\) 层分配了一个预测中心，但每个层内部的数据分布还在波动。总方差应该可以写成层内波动与层间波动之和。

从恒等式出发：

\[
X-\E[X]=\underbrace{(X-\E[X\mid Y])}_{\text{层内残差}}+\underbrace{(\E[X\mid Y]-\E[X])}_{\text{层间偏离}}.
\]

两边平方再取期望。交叉项是

\[
2\E\bigl[(X-\E[X\mid Y])(\E[X\mid Y]-\E[X])\bigr].
\]

考察它：先对 \(Y\) 取条件期望。给定 \(Y\) 时，\(\E[X\mid Y]-\E[X]\) 由 \(Y\) 的值完全确定，是个常数，可以提到条件期望外：

\[
\E\bigl[(X-\E[X\mid Y])(\E[X\mid Y]-\E[X])\mid Y\bigr]
=(\E[X\mid Y]-\E[X])\,\E\bigl[X-\E[X\mid Y]\mid Y\bigr].
\]

而 \(\E[X-\E[X\mid Y]\mid Y]=\E[X\mid Y]-\E[X\mid Y]=0\)。给定 \(Y\) 后，残差的期望精确为零——这是条件期望的定义性质。所以整个交叉项的期望为零，得到

\[
\Var(X)=\underbrace{\E[\Var(X\mid Y)]}_{\text{层内波动}}+\underbrace{\Var(\E[X\mid Y])}_{\text{层间波动}}.
\]

两项的含义：

\begin{itemize}
\item \textbf{层内波动} \(\E[\Var(X\mid Y)]\)：即使告诉你温度，每一天用电仍围绕当天温度的典型值在浮动。这是信息 \(Y\) 也解释不了的那部分随机性。
\item \textbf{层间波动} \(\Var(\E[X\mid Y])\)：不同温度条件下预测均值本身的差异。这是信息 \(Y\) 能捕获的那部分总体变异。
\end{itemize}

用电例子配数字。假设高温天的条件方差是 15000，低温天是 9000。则

\[
\E[\Var(X\mid Y)]=0.3\times15000+0.7\times9000=4500+6300=10800.
\]

不同温度预测中心的方差：

\[
\Var(\E[X\mid Y])=0.3\times(650-500)^2+0.7\times(436-500)^2=0.3\times22500+0.7\times4096=6750+2867=9617.
\]

总方差 \(\Var(X)=10800+9617=20417\)。温度信息解释了约 \(9617/20417\approx 47\%\) 的总变异。在\hyperref[sec:05-Probability/02-Conditional-Probability/02]{``全概率公式与 Bayes 更新''}的机器次品例子中，同样的分解告诉你次品率的总波动有多少来自机器间的差异、多少来自同一台机器内部的随机次品——前者可以靠改进设备降低，后者或许需要改进工艺。

这个分解在实际数据分析中就是 ANOVA 变差分解的概率版，也是理解 ``某个特征有多大解释力'' 的基础框架。

\subsection{条件期望是所有 \(Y\) 的函数中预测最准的那个}

现在回到开头的判断问题：用什么样的 \(g(Y)\) 预测 \(X\) 使得 \(\E[(X-g(Y))^2]\) 最小？

把预测误差拆成两段，插入 \(\E[X\mid Y]\) 作为参照点：

\[
\begin{aligned}
\E[(X-g(Y))^2]&=\E[(X-\E[X\mid Y]+\E[X\mid Y]-g(Y))^2]\\
&=\E[(X-\E[X\mid Y])^2]+\E[(\E[X\mid Y]-g(Y))^2]\\
&\quad+2\E[(X-\E[X\mid Y])(\E[X\mid Y]-g(Y))].
\end{aligned}
\]

交叉项的处理方式与全方差公式一致：给定 \(Y\) 时 \(\E[X\mid Y]-g(Y)\) 由 \(Y\) 确定，残差的条件期望为零，交叉项消失。

第一项 \(\E[(X-\E[X\mid Y])^2]\) 与你的 \(g\) 选择无关——它是条件期望本身固有的残余误差，你改不了。第二项 \(\E[(\E[X\mid Y]-g(Y))^2]\) 衡量你的 \(g\) 与条件期望之间的差距，总是非负，且只有 \(g(Y)=\E[X\mid Y]\) 时归零。因此：

\[
\min_{g}\E[(X-g(Y))^2]=\E[(X-\E[X\mid Y])^2],
\]

最优解就是 \(g^*(Y)=\E[X\mid Y]\)。条件期望是使用信息 \(Y\) 时，均方误差意义下的最优预测器。

这组论证还有几何面向。把随机变量看作内积 \(\E[UV]\) 下的 Hilbert 空间的 ``向量''。条件期望 \(\E[X\mid Y]\) 是 \(X\) 向 ``可由 \(Y\) 构造的函数'' 这一子空间的正交投影。残差 \(X-\E[X\mid Y]\) 与该子空间内所有向量正交：

\[
\E[(X-\E[X\mid Y])\,h(Y)]=0,
\]

对任意使期望有意义的 \(h\) 成立。验证：先对 \(Y\) 取条件期望，\(h(Y)\) 是已知量可提出，剩下 \(\E[X-\E[X\mid Y]\mid Y]=0\)。这一结构与最小二乘回归、Kalman 滤波和 Hilbert 空间投影理论全部连通。

\subsection{几个它不负责的事}

和任何强大的概念一样，条件期望的局限恰恰是你在用之前应该知道的东西。

\textbf{最优不等于线性。} \(\E[X\mid Y]\) 一般是 \(Y\) 的非线性函数。如果你只允许形如 \(aY+b\) 的预测器，得到的是最佳线性预测——不等于真正的条件期望。两者的差距衡量了依赖关系中 ``线性部分之外'' 有多少。联合 Gaussian 是少有的例外族：在其中条件期望对条件变量是仿射的。

\textbf{连续变量时不能直接用事件概率做除法。} 如果 \(Y\) 是连续变量（比如温度取值连续），对每个精确值 \(Y=y\)，\(P(Y=y)=0\)。``条件概率'' \(P(X\le x\mid Y=y)\) 不再能用两个概率的比值来定义，因为分母为零。这需要更精细的工具——条件期望的形式定义通过可测性和局部积分匹配来绕开这个困难。\hyperref[ch:101]{``测度、积分与条件期望：为现代概率回填地基''}会专门处理这一点，并说明为什么不同版本的 \(\E[X\mid Y]\) 只需要几乎处处相等。

\textbf{最佳仅限于平方损失。} 我们的全部论述都建立在 \(\E[(X-g(Y))^2]\) 的极小化上。如果你关心的是绝对误差 \(\E[|X-g(Y)|]\)，最优预测是条件中位数而不是条件期望。如果损失函数左右不对称——比如低估需求的惩罚远高于高估——最优预测可能是某个条件分位数。换了损失函数就得换最优预测器，没有一成不变的 ``最佳''。

\textbf{条件期望只定义了几乎处处。} 在测度论框架中，任何两个函数只要在不大于一个零测集上有差别，都被视为 \(\E[X\mid Y]\) 的有效版本。实践中这不碍事——改变零概率事件上的值不影响任何概率计算。但理论上你永远在处理等价类，不是唯一函数。

这些不是缺陷，是边界。知道条件期望在自家领地内是最优预测器，同时也知道领地的围墙在哪里，才能正确地走下一步——走到最优线性预测、走到非线性回归、走到条件中位数与分位数回归，或者走到\hyperref[ch:041]{``联合分布：单个分布之外，变量怎样一起变化''}，用联合分布框架统一处理多个变量同时变化的全部信息。
